\documentclass[../../main.tex]{subfiles}

\graphicspath{{../../image/}{../../figure/}}

\begin{document}

\section{对合与偶然熵}\label{section:ch05-对合与偶然熵}

前三节展示了凸熵的有益作用。然而，连续介质物理中平衡律方程组所伴随的熵并不总是凸的。事实上，我们已在第 III 章中遇到过等熵弹性动力学（\cref{section:ch03-例子}）与电动力学（\cref{section:ch03-例子}）的情形，

其中由物理决定的不变性与熵的全局凸性不相容。本节及下一节的目标是识别此类方程组中可补偿熵凸性缺失的特殊结构。

回忆与弹性动力学相关的方程组\eqref{align:ch03-等熵热弹性系统}的解还应满足方程\eqref{align:ch03-协调方程}。

注意\eqref{align:ch03-协调方程}并不独立于\eqref{align:ch03-等熵热弹性系统}。

事实上，在 Cauchy 问题中，\eqref{align:ch03-等熵热弹性系统}$_{1}$
蕴涵当\eqref{align:ch03-协调方程}被初值满足时，它对所有 $t>0$ 成立。

电动力学方程组表现出类似行为：除双曲方程组\eqref{align:ch03-Maxwell}外，
磁感应与电位移还须满足\eqref{align:ch03-高斯}。
然而在 Cauchy 问题中，由\eqref{align:ch03-Maxwell}与\eqref{align:ch03-电荷守恒}，

只要\eqref{align:ch03-高斯}中两个等式被初值满足，它们对所有 $t>0$ 自动成立。

人们在连续介质物理中出现的许多其他方程组中识别出类似结构，因此有必要在一般框架下考察其蕴含。

我们考虑双曲方程组\eqref{align:ch05-cauchy-system}中具有如下对称性条件性质的一类：
\begin{align}
M_{\alpha}G_{\beta}(U)+M_{\beta}G_{\alpha}(U)=0,\qquad U\in\overline{\mathscr{B}_{\rho}},\ \alpha,\beta=1,\ldots,m,\label{align:ch05-symmetry-condition}
\end{align}
对某 $k\times n$ 矩阵族 $M_{\alpha}$、$\alpha=1,\ldots,m$ 成立。\eqref{align:ch05-symmetry-condition}的直接推论是：\eqref{align:ch05-cauchy-system} Cauchy 问题的任一（一般为弱）解 $U$ 将满足附加方程
\begin{align}
\sum_{\alpha=1}^{m}M_{\alpha}\partial_{\alpha}U=0,\label{align:ch05-involution}
\end{align}
只要初值满足它。我们称\eqref{align:ch05-involution}为\eqref{align:ch05-cauchy-system}的一个\emph{对合}（involution）。于是\eqref{align:ch03-协调方程}是\eqref{align:ch03-等熵热弹性系统}与\eqref{align:ch03-热弹性系统}共同的对合，


而\eqref{align:ch03-高斯}是\eqref{align:ch03-Maxwell}的对合。

典型地，对物理中出现的此类方程组，唯一相关的解是同时满足对合的解。

对任意 $v\in\mathbb{S}^{m-1}$，我们伴随 $k\times n$ 矩阵
\begin{align}
N(v)=\sum_{\alpha=1}^{m}v_{\alpha}M_{\alpha}.\label{align:ch05-Nv-def}
\end{align}
由\eqref{align:ch04-方向矩阵}与\eqref{align:ch05-symmetry-condition}，
\begin{align}
N(v)\Lambda(v;U)=0,\qquad v\in\mathbb{S}^{m-1},\ U\in\overline{\mathscr{B}_{\rho}},\label{align:ch05-N-Lambda}
\end{align}
这表明在对合存在时，$0$ 必为 $\Lambda(v;U)$ 的特征值，其几何重数至少等于 $N(v)$ 的秩；且 $\Lambda(v;U)$ 的任一非零特征值 $\lambda_{i}(v;U)$ 对应的特征向量 $R_{i}(v;U)$ 必落在 $N(v)$ 的核中。

还应注意，与满足对合\eqref{align:ch05-involution}的\eqref{align:ch05-cauchy-system}解相联系的、沿方向 $v$ 传播的任一激波必满足跳跃条件
\begin{align}
N(v)[U_{+}-U_{-}]=0,\label{align:ch05-shock-jump}
\end{align}
从而其振幅 $U_{+}-U_{-}$ 必落在 $N(v)$ 的核中。

本节中我们将在如下假设下工作：
\begin{align}
\dim\ker\Lambda(v;U)=\operatorname{rank}N(v),\qquad v\in\mathbb{S}^{m-1},\ U\in\overline{\mathscr{B}_{\rho}},\label{align:ch05-kernel-rank}
\end{align}
此时 $N(v)$ 的核由 $\Lambda(v;U)$ 的非零特征值对应的特征向量张成，从而与 $\Lambda(v;U)$ 的值域重合。对等熵弹性动力学方程组\eqref{align:ch03-等熵热弹性系统}，
情况正是如此，其中 $\Lambda(v;U)$ 的核为六维而 $N(v)$ 的秩为六；但对非等熵热弹性方程组\eqref{align:ch03-热弹性系统}则不然，
其中 $\Lambda(v;U)$ 的核为七维。带对合的方程组（如\eqref{align:ch03-热弹性系统}，其中 $\ker\Lambda(v;U)$ 的维数大于 $N(v)$ 的秩）可由同一方法处理，代价是引入附加结构——读者应免受此类复杂情况之扰。

现引入\emph{对合锥}
\begin{align}
\mathscr{C}=\bigcup_{v\in\mathbb{S}^{m-1}}\ker N(v),\label{align:ch05-involution-cone}
\end{align}
嵌入状态空间 $\mathbb{R}^{n}$ 中。特别地，对等熵弹性动力学方程组\eqref{align:ch03-等熵热弹性系统}，
$\mathscr{C}=\{(F,v):F=uw^{\top},\ u,w,v\in\mathbb{R}^{3}\}$；而对电动力学方程组\eqref{align:ch03-Maxwell}，
$\mathscr{C}$ 占据整个状态空间 $\mathbb{R}^{6}$。正如我们将看到的，一切有趣的活动都发生在 $\mathscr{C}$ 上。

对合的存在使熵概念自然推广。回忆\eqref{align:ch05-cauchy-system}的熵--熵通量对 $(\eta,Q)$ 的定义被构造为使额外守恒律
\begin{align}
\partial_{t}\eta(U(x,t))+\sum_{\alpha=1}^{m}\partial_{\alpha}Q_{\alpha}(U(x,t))=0\label{align:ch05-entropy-conservation}
\end{align}
自动被\eqref{align:ch05-cauchy-system}的任一 $C^{1}$ 解满足。然而在当前设定下，合理的要求是\eqref{align:ch05-entropy-conservation}仅对同时满足对合\eqref{align:ch05-involution}的 $C^{1}$ 解恒等成立。这引出如下

\begin{definition}[偶然熵]\label{definition:ch05-contingent-entropy}
在装备有对合\eqref{align:ch05-involution}的守恒律方程组\eqref{align:ch05-cauchy-system}中，$\overline{\mathscr{B}_{\rho}}$ 上光滑的标量值函数 $\eta$ 连同 $1\times k$ 矩阵值的\emph{偶然熵通量} $Q(U)$ 一起称为\emph{偶然熵}（contingent entropy），若存在 $\overline{\mathscr{B}_{\rho}}$ 上 $k$-向量值函数 $\Xi(U)$ 使得
\begin{align}
\mathrm{D}Q_{\alpha}(U)=\mathrm{D}\eta(U)\mathrm{D}G_{\alpha}(U)+\Xi(U)^{\top}M_{\alpha},\qquad\alpha=1,\ldots,m.\label{align:ch05-contingent-def}
\end{align}
\end{definition}

特别地，任一熵都是偶然熵，取 $\Xi=0$。另一方面，由\eqref{align:ch03-余因子守恒}与\eqref{align:ch03-行列式守恒}，

$\det F$ 与矩阵 $F^{*}$ 的九个元素都是等熵弹性动力学方程组\eqref{align:ch03-等熵热弹性系统}的偶然熵，但不是熵。
类似地，由\eqref{align:ch03-动量电磁}，
$B\wedge D$ 的三个分量是电动力学方程组\eqref{align:ch03-Maxwell}的偶然熵，但不是熵。
这些特定偶然熵的有用作用将在下一节阐明。\eqref{align:ch03-等熵热弹性系统}的另一个有趣的偶然熵例子在此处不会使用，
即 3-向量 $F^{\top}v$ 的各分量。相应的偶然熵通量是 $3\times3$ 矩阵 $F^{\top}S+(\varepsilon+\frac{1}{2}|v|^{2})I$ 的相应行。

\eqref{align:ch05-contingent-def}的可积性条件推广了对称关系\eqref{align:ch03-可积性}，
即 $n\times n$ 矩阵
\begin{align}
J_{\alpha}(U)=A(U)\mathrm{D}G_{\alpha}(U)+\mathrm{D}\Xi(U)^{\top}M_{\alpha},\qquad\alpha=1,\ldots,m,\label{align:ch05-Jalpha-contingent}
\end{align}
对称。

\cref{section:ch05-凸熵与经典解的存在性}的结果不能应用于等熵弹性动力学方程组\eqref{align:ch03-等熵热弹性系统}，

因为如我们所见，熵 $\eta=\varepsilon(F)+\frac{1}{2}|v|^{2}$ 的全局凸性与物理定律不相容。然而，物理上合理的假设\eqref{align:ch03-Legendre-Hadamard}——
即内能 $\varepsilon(F)$ 秩一凸——保证了 $\eta$ 至少在对合锥 $\mathscr{C}$ 上凸。本节的目的是证明：在对合\eqref{align:ch05-involution}存在时，在方程组\eqref{align:ch05-cauchy-system}装备有仅在对合锥上凸的偶然熵 $\eta$ 的假设下，可建立 Cauchy 问题经典解的局部存在性与稳定性，即
\begin{align}
X^{\top}A(U)X\geqslant 2\mu|X|^{2},\qquad X\in\mathscr{C},\ U\in\overline{\mathscr{B}_{\rho}},\label{align:ch05-cone-convex}
\end{align}
其中 $\mu>0$。

取代\cref{theorem:ch05-凸熵局部存在}，我们此处有如下命题：

\begin{theorem}\label{theorem:ch05-involution-existence}
假设守恒律双曲方程组\eqref{align:ch05-cauchy-system}满足\eqref{align:ch05-symmetry-condition}、\eqref{align:ch05-kernel-rank}，并装备有在对合锥 $\mathscr{C}$ 上凸的偶然熵 $\eta$，即\eqref{align:ch05-cone-convex}成立。设初值 $U_{0}\in H_{\ell}$，某 $\ell>\frac{m}{2}+1$，取值于半径 $\rho_{0}<\rho$ 的球 $\overline{\mathscr{B}_{\rho_{0}}}$，并满足对合\eqref{align:ch05-involution}。则存在 $T_{\infty}\leqslant\infty$ 及 $\mathbb{R}^{m}\times[0,T_{\infty})$ 上唯一连续可微、取值于 $\overline{\mathscr{B}_{\rho}}$ 的函数 $U$，它是 Cauchy 问题\eqref{align:ch05-cauchy-system}、\eqref{align:ch05-cauchy-data}在时间区间 $[0,T_{\infty})$ 上的经典解。此外，
\begin{align}
U(\cdot,t)\in\bigcap_{k=0}^{\ell}C^{k}([0,T_{\infty});H_{\ell-k}).\label{align:ch05-involution-regularity}
\end{align}
区间 $[0,T_{\infty})$ 是极大的，即若 $T_{\infty}<\infty$，则
\begin{align}
\int_{0}^{T_{\infty}}\|\|\nabla U(\cdot,t)\|\|_{L^{\infty}}\,\mathrm{d}t=\infty,\label{align:ch05-involution-blowup1}
\end{align}
且/或
\begin{align}
\limsup_{t\to T_{\infty}}\|U(\cdot,t)\|_{L^{\infty}}=\rho.\label{align:ch05-involution-blowup2}
\end{align}
\end{theorem}

\begin{theorem}\label{theorem:ch05-involution-stability}
假设守恒律双曲方程组\eqref{align:ch05-cauchy-system}满足\eqref{align:ch05-symmetry-condition}、\eqref{align:ch05-kernel-rank}，并装备有偶然熵--熵通量对 $(\eta,Q)$，其中 $\eta$ 在对合锥 $\mathscr{C}$ 上凸，即\eqref{align:ch05-cone-convex}成立。设 $\bar{U}$ 是\eqref{align:ch05-cauchy-system}在 $[0,T)$ 上的经典解，取值于 $\overline{\mathscr{B}_{\rho}}$，初值 $\bar{U}_{0}$ 满足对合\eqref{align:ch05-involution}。此外，设 $\bar{U}(x,t)\to0$，$|x|\to\infty$，对 $t\in[0,T)$ 一致。设 $U$ 是\eqref{align:ch05-cauchy-system}在 $[0,T)$ 上的任一弱解，取值于 $\overline{\mathscr{B}_{\rho}}$，满足熵可容许性条件\eqref{align:ch04-熵积分0}，
且初值 $U_{0}$ 满足对合\eqref{align:ch05-involution}。此外，设 $U(x,t)\to0$，$|x|\to\infty$，对 $t\in[0,T)$ 一致。最后，设 $\|U(\cdot,t)-\bar{U}(\cdot,t)\|_{0}$ 在 $[0,T)$ 上有界。在上述假设下，存在常数 $a$ 与 $\kappa$，依赖于 $\rho$ 以及 $G,\eta,Q$ 及其在 $\overline{\mathscr{B}_{\rho}}$ 上导数的界，以及还依赖于 $\bar{U}$ Lipschitz 常数的 $b$，使得若
\begin{align}
\limsup_{x\to y,\ t\to\tau}|U(x,t)-U(y,\tau)|<\kappa,\qquad y\in\mathbb{R}^{m},\ \tau\in[0,T),\label{align:ch05-small-oscillation}
\end{align}
则
\begin{align}
\int_{\mathbb{R}^{m}}|U(x,t)-\bar{U}(x,t)|^{2}\,\mathrm{d}x\leqslant ae^{bt}\int_{\mathbb{R}^{m}}|U_{0}(x)-\bar{U}_{0}(x)|^{2}\,\mathrm{d}x,\qquad a.e.\ \text{在 }[0,T)\text{ 上}.\label{align:ch05-involution-stability-est}
\end{align}
特别地，$\bar{U}$ 是 Cauchy 问题\eqref{align:ch05-cauchy-system}在初值 $\bar{U}_{0}$、具有充分小局部振荡\eqref{align:ch05-small-oscillation}且与 $\bar{U}$ 在 $|x|=\infty$ 处具有相同渐近行为的可容许弱解类中的唯一解。
\end{theorem}

下述引理揭示了对合如何补偿熵在对合锥外的凸性缺失，将在上述两个命题的证明中起关键作用。

\begin{lemma}\label{lemma:ch05-compactness}
设 $P$ 是 $\mathbb{R}^{m}$ 上有界可测的对称 $n\times n$ 矩阵值函数，满足
\begin{align}
X^{\top}P(x)X\geqslant 2\mu|X|^{2},\qquad X\in\mathscr{C},\ x\in\mathbb{R}^{m}.\label{align:ch05-P-positive}
\end{align}
进一步假设 $\mathbb{R}^{m}$ 存在由开集 $\Omega_{0},\Omega_{1},\ldots,\Omega_{K}$ 的并组成的有限覆盖，具有性质：对 $J=0,1,\ldots,K$，
\begin{align}
|P(x)-P(y)|<\frac{1}{2}\mu,\qquad x,y\in\Omega_{J}.\label{align:ch05-P-oscillation}
\end{align}
则存在 $\delta$，仅依赖于该覆盖，使得对任一满足对合
\begin{align}
\sum_{\alpha=1}^{m}M_{\alpha}\partial_{\alpha}S=0,\label{align:ch05-S-involution}
\end{align}
（在 $\mathbb{R}^{m}$ 上分布意义下）的 $S\in L^{2}(\mathbb{R}^{m};\mathbb{R}^{n})$，成立
\begin{align}
\int_{\mathbb{R}_{m}}S(x)^{\top}P(x)S(x)\,\mathrm{d}x\geqslant\mu\|S\|_{0}^{2}-\delta\|S\|_{-1}^{2}.\label{align:ch05-compactness-est}
\end{align}
\end{lemma}

\begin{proof}
固定 $U\in\overline{\mathscr{B}_{\rho}}$，考虑线性微分算子
\begin{align}
\mathscr{L}=\sum_{\alpha=1}^{m}\mathrm{D}G_{\alpha}(U)\partial_{\alpha}.\label{align:ch05-L-def}
\end{align}
我们借助 Fourier 变换求方程
\begin{align}
\mathscr{L}\Phi+\Phi=S,\label{align:ch05-L-Phi}
\end{align}
的解 $\Phi\in H_{1}$：
\begin{align}
[i|\xi|\Lambda(v;U)+I]\hat{\Phi}(\xi)=\hat{S}(\xi),\qquad\xi\in\mathbb{R}^{m},\label{align:ch05-fourier}
\end{align}
其中 $v=|\xi|^{-1}\xi$，$\Lambda(v;U)$ 由\eqref{align:ch04-方向矩阵}定义。
\eqref{align:ch05-fourier}左端的 $n\times n$ 矩阵非奇异，因为 $\Lambda(v;U)$ 仅有实特征值。此外，由\eqref{align:ch05-involution}，$\hat{S}(\xi)$ 落在 $N(v)$ 的核中，而后者由 $\Lambda(v;U)$ 的非零特征值对应的特征向量张成。于是\eqref{align:ch05-fourier}有解 $\hat{\Phi}(\xi)$ 满足
\begin{align}
|\hat{\Phi}(\xi)|^{2}\leqslant c^{2}(1+|\xi|^{2})^{-1}|\hat{S}(\xi)|^{2},\qquad\xi\in\mathbb{R}^{m},\label{align:ch05-Phi-bound}
\end{align}
从而
\begin{align}
\|\Phi\|_{0}\leqslant c\|S\|^{-1}.\label{align:ch05-Phi-norm}
\end{align}

接下来考虑单位分解 $\psi_{0},\psi_{1},\ldots,\psi_{K}$，从属于覆盖 $\Omega_{0},\Omega_{1},\ldots,\Omega_{K}$，即对 $J=0,1,\ldots,K$，$\psi_{J}\in C^{\infty}(\mathbb{R}^{m})$，$\operatorname{spt}\psi_{J}\subset\Omega_{J}$，且
\begin{align}
\sum_{J=0}^{K}\psi_{J}^{2}(x)=1,\qquad x\in\mathbb{R}^{m}.\label{align:ch05-partition}
\end{align}
再固定 $y_{J}\in\Omega_{J}$，$J=0,1,\ldots,K$，写
\begin{align}
\int_{\mathbb{R}^{m}}S(x)^{\top}P(x)S(x)\,\mathrm{d}x
&=\sum_{J=0}^{K}\int_{\mathbb{R}^{m}}\psi_{J}^{2}(x)S(x)^{\top}P(x)S(x)\,\mathrm{d}x\nonumber\\
&=\sum_{J=0}^{K}\int_{\mathbb{R}^{m}}\psi_{J}^{2}(x)S(x)^{\top}P(y_{J})S(x)\,\mathrm{d}x\nonumber\\
&\quad+\sum_{J=0}^{K}\int_{\mathbb{R}^{m}}\psi_{J}^{2}(x)S(x)^{\top}[P(x)-P(y_{J})]S(x)\,\mathrm{d}x.\label{align:ch05-P-split}
\end{align}
由\eqref{align:ch05-P-oscillation}，
\begin{align}
\sum_{J=0}^{K}\int_{\mathbb{R}^{m}}\psi_{J}^{2}(x)S(x)^{\top}[P(x)-P(y_{J})]S(x)\,\mathrm{d}x\geqslant-\frac{1}{2}\mu\|S\|_{0}^{2}.\label{align:ch05-P-oscillation-est}
\end{align}
对每个 $J=0,1,\ldots,K$，我们将 $\psi_{J}S$ 分裂为
\begin{align}
\psi_{J}S=X_{J}+Y_{J},\label{align:ch05-XY-split}
\end{align}
其中
\begin{align}
X_{J}=\mathscr{L}(\psi_{J}\Phi),\label{align:ch05-XJ-def}
\end{align}
\begin{align}
Y_{J}=\left[\psi_{J}I-\sum_{\alpha=1}^{m}\partial_{\alpha}\psi_{J}\mathrm{D}G_{\alpha}(U)\right]\Phi.\label{align:ch05-YJ-def}
\end{align}
置 $v=|\xi|^{-1}\xi$ 并回忆\eqref{align:ch05-N-Lambda}，
\begin{align}
N(v)\hat{X}_{J}(\xi)=i|\xi|N(v)\Lambda(v,U)(\widehat{\psi_{J}\Phi})(\xi)=0,\label{align:ch05-XJ-kernel}
\end{align}
于是对任意 $\xi\in\mathbb{R}^{m}$ 及 $J=0,1,\ldots,K$，$\hat{X}_{J}(\xi)$ 的实部与虚部都落在 $\mathscr{C}$ 中。因此，应用 Parseval 关系并利用\eqref{align:ch05-P-positive}，我们推得
\begin{align}
\int_{\mathbb{R}^{m}}X_{J}(x)^{\top}P(y_{J})X_{J}(x)\,\mathrm{d}x=\int_{\mathbb{R}^{m}}\hat{X}_{J}(\xi)^{*}P(y_{J})\hat{X}_{J}(\xi)\,\mathrm{d}\xi\geqslant 2\mu\int_{\mathbb{R}^{m}}|X_{J}(x)|^{2}\,\mathrm{d}x.\label{align:ch05-XJ-positive}
\end{align}
此外，由\eqref{align:ch05-YJ-def}与\eqref{align:ch05-Phi-norm}，我们推得
\begin{align}
\int_{\mathbb{R}^{m}}|Y_{J}(x)|^{2}\,\mathrm{d}x\leqslant c_{1}\|S\|_{-1}^{2},\qquad J=0,\ldots,K.\label{align:ch05-YJ-bound}
\end{align}

回到\eqref{align:ch05-P-split}。由\eqref{align:ch05-XY-split}、\eqref{align:ch05-XJ-positive}与\eqref{align:ch05-YJ-bound}，
\begin{align}
\int_{\mathbb{R}^{m}}\psi_{J}^{2}(x)S(x)^{\top}P(y_{J})S(x)\,\mathrm{d}x
&\geqslant\frac{7}{8}\int_{\mathbb{R}^{m}}X_{J}(x)^{\top}P(y_{J})X_{J}(x)\,\mathrm{d}x-8\int_{\mathbb{R}^{m}}Y_{J}(x)^{\top}P(y_{J})Y_{J}(x)\,\mathrm{d}x\nonumber\\
&\geqslant\frac{7}{4}\mu\int_{\mathbb{R}^{m}}|X_{J}(x)|^{2}\,\mathrm{d}x-c_{2}\|S\|_{-1}^{2}.\label{align:ch05-PJ-est}
\end{align}
又由\eqref{align:ch05-XY-split}与\eqref{align:ch05-YJ-bound}，
\begin{align}
\int_{\mathbb{R}^{m}}|X_{J}(x)|^{2}\,\mathrm{d}x\geqslant\frac{6}{7}\int_{\mathbb{R}^{m}}\psi_{J}^{2}(x)|S(x)|^{2}\,\mathrm{d}x-c_{3}\|S\|_{-1}^{2}.\label{align:ch05-XJ-lower}
\end{align}
合并\eqref{align:ch05-partition}、\eqref{align:ch05-P-split}、\eqref{align:ch05-P-oscillation-est}、\eqref{align:ch05-PJ-est}与\eqref{align:ch05-XJ-lower}，即得\eqref{align:ch05-compactness-est}。引理证明完成。
\end{proof}

\begin{proof}
\emph{\cref{theorem:ch05-involution-existence}的证明。}在当前假设下，\eqref{align:ch05-cauchy-system}、\eqref{align:ch05-cauchy-data}的解将由粘性消去法构造，即作为同初值\eqref{align:ch05-cauchy-data}的抛物方程组\eqref{align:ch05-parabolic}的解当 $\varepsilon\to0$ 时的极限。这一在\cref{section:ch05-凸熵与经典解的存在性}已检验过的途径在当前设定下同样有效，
因为当\eqref{align:ch05-symmetry-condition}成立时，抛物方程组\eqref{align:ch05-parabolic}从双曲方程组\eqref{align:ch05-cauchy-system}继承了对合\eqref{align:ch05-involution}。\footnote{注意线性化方程组\eqref{align:ch05-linearized}并非如此，因此在对合存在时，传统方法在当前假设下似乎不适用。}事实上，若 $U$ 是\eqref{align:ch05-parabolic}以满足对合\eqref{align:ch05-involution}的初值\eqref{align:ch05-cauchy-data}的解，则 $Z=\sum M_{\alpha}\partial_{\alpha}U$ 是热方程 $\partial_{t}Z=\varepsilon\Delta Z$ 以零初值为初值的解，从而恒等于零。

当前设定下解的构造将紧密平行于\cref{theorem:ch05-凸熵局部存在}中经典情形的处理。
由于 $\eta$ 不再是凸熵而仅为在对合锥上凸的偶然熵，需要若干修改。为最小化重复，我们不写冗长的自含证明，而仅回溯\cref{theorem:ch05-凸熵局部存在}证明中经由引理 5.1.2、5.1.3、5.1.4 的步骤，视需要插入调整。

首先，\cref{lemma:ch05-parabolic-local}并不依赖于熵的存在，因此此处无需修改即适用。于是我们知道，对固定 $\varepsilon>0$ 与 $\omega>\|U_{0}\|_{\ell}$，存在\eqref{align:ch05-parabolic}、\eqref{align:ch05-cauchy-data}在时间区间 $[0,T_{\omega,\varepsilon})$ 上的解 $U$，取值于 $\overline{\mathscr{B}_{\rho}}$，满足\eqref{align:ch05-lemma-bound}。我们回溯\cref{lemma:ch05-uniform-bound}证明的步骤，导出 $U$ 不依赖于 $\varepsilon$ 的界。

\eqref{align:ch05-Ur-eq2}、\eqref{align:ch05-Moser}、\eqref{align:ch05-id1}、\eqref{align:ch05-id2}、\eqref{align:ch05-id3}与\eqref{align:ch05-id4}仍然成立。我们不再能依靠矩阵 $A(U)\mathrm{D}G_{\alpha}(U)$ 的对称性。然而，当改用\eqref{align:ch05-Jalpha-contingent}作为 $J_{\alpha}(U)$ 的定义时，\eqref{align:ch05-id5}仍然成立，因为 $U$ 从而 $U_{r}$ 满足对合\eqref{align:ch05-involution}。作此修改后，基本"能量"积分\eqref{align:ch05-energy-A}仍然有效。

下一个障碍是我们不再有\eqref{align:ch05-A-positive}与\eqref{align:ch05-sum-positive}，因为 $A(U)$ 不必正定。我们将借助\cref{lemma:ch05-compactness}补偿 $\eta(U)$ 凸性的缺失。由于 $U_{0}\in H_{\ell}$，
\begin{align}
|A(U_{0}(x))-A(U_{0}(y))|<\frac{1}{4}\mu,\label{align:ch05-A0-oscillation}
\end{align}
对集合 $\Omega_{0}=\{z\in\mathbb{R}^{m}:|z|>\alpha\}$ 中所有 $x,y$ 成立，其中 $\alpha$ 充分大。接着以半径充分小的球 $\Omega_{1},\ldots,\Omega_{K}$ 覆盖紧集 $\Omega_{0}^{c}$，使得\eqref{align:ch05-A0-oscillation}对任意 $x,y\in\Omega_{I}$、$I=1,\ldots,K$ 亦成立。最后回忆\eqref{align:ch05-Linf-bound}，将 $U$ 限制到时间区间 $[0,T]$，$T<T_{\omega,\varepsilon}$，使得
\begin{align}
|A(U(x,t))-A(U_{0}(x))|<\frac{1}{8}\mu,\label{align:ch05-AU-oscillation}
\end{align}
对所有 $x\in\mathbb{R}^{m}$ 与 $t\in[0,T]$ 成立。由此，对任意固定 $t\in[0,T]$，$\mathbb{R}^{m}$ 的覆盖 $\Omega_{0},\Omega_{1},\ldots,\Omega_{K}$ 满足\cref{lemma:ch05-compactness}中的条件，以 $A(U(\cdot,t))$ 扮演 $P(\cdot)$。特别地，取代\eqref{align:ch05-A-positive}与\eqref{align:ch05-sum-positive}，我们现已有
\begin{align}
\int_{\mathbb{R}^{m}}\partial_{\alpha}U_{r}^{\top}(x,\tau)A(U(x,\tau))\partial_{\alpha}U_{r}(x,\tau)\,\mathrm{d}x
\geqslant\mu\int_{\mathbb{R}^{m}}|\partial_{\alpha}U_{r}(x,\tau)|^{2}\,\mathrm{d}x-\delta\int_{\mathbb{R}^{m}}|U_{r}(x,\tau)|^{2}\,\mathrm{d}x,\label{align:ch05-A-cones}
\end{align}
\begin{align}
\sum_{r\leqslant\ell}\int_{\mathbb{R}^{m}}U_{r}^{\top}(x,t)A(U(x,t))U_{r}(x,t)\,\mathrm{d}x\geqslant\mu\|U(\cdot,t)\|_{\ell}^{2}-c\|U(\cdot,t)\|_{\ell-1}^{2}.\label{align:ch05-sum-cones}
\end{align}

因此，\eqref{align:ch05-A-cones}连同\eqref{align:ch05-energy-A}、\eqref{align:ch05-CS}与\eqref{align:ch05-Moser}表明\eqref{align:ch05-sum-energy}在此处仍然有效，其中
\begin{align}
g(\tau)\leqslant c[1+\|\|\nabla U(\cdot,\tau)\|\|_{L^{\infty}}+\varepsilon\|\|\nabla U(\cdot,\tau)\|\|_{L^{\infty}}^{2}].\label{align:ch05-g-cones}
\end{align}

参照\eqref{align:ch05-sum-cones}，为估计项 $\|U(\cdot,t)\|_{\ell-1}$，我们以 $2U_{r}^{\top}$ 乘\eqref{align:ch05-Ur-eq2}，在 $\mathbb{R}^{m}\times[0,t]$、$t\in[0,T]$ 上积分，分部积分项 $2\varepsilon U_{r}^{\top}\Delta U_{r}$，并对所有阶 $|r|\leqslant\ell-1$ 的多重指标 $r$ 求和。利用\eqref{align:ch05-Moser}，我们推得
\begin{align}
\|U(\cdot,t)\|_{\ell-1}^{2}\leqslant\|U_{0}(\cdot)\|_{\ell-1}^{2}+c\int_{0}^{t}[1+\|\|\nabla U(\cdot,\tau)\|\|_{L^{\infty}}]\|U(\cdot,\tau)\|_{\ell}^{2}\,\mathrm{d}\tau.\label{align:ch05-Ul1-est}
\end{align}

合并\eqref{align:ch05-sum-energy}与\eqref{align:ch05-sum-cones}、\eqref{align:ch05-Ul1-est}，我们推出\eqref{align:ch05-Gronwall-setup}、从而\eqref{align:ch05-uniform-est}在当前设定下成立，$g$ 通过\eqref{align:ch05-g-cones}有界。

如\cref{lemma:ch05-uniform-bound}中那样，我们结论：对任意固定 $\omega>c_{0}\|U_{0}\|_{\ell}$，满足 $\|U(\cdot,t)\|_{\ell}<\omega$ 的\eqref{align:ch05-parabolic}、\eqref{align:ch05-cauchy-data}的解 $U$ 对任意 $\varepsilon>0$ 在时间区间 $[0,T_{\omega})$ 上存在。

回溯\cref{lemma:ch05-limit}证明的步骤，我们现将 Cauchy 问题\eqref{align:ch05-cauchy-system}、\eqref{align:ch05-cauchy-data}的解 $U$ 构造为解序列 $\{U_{k}\}$ 当 $\varepsilon=\varepsilon_{k}$、$\varepsilon_{k}\to0$（$k\to\infty$）时的极限。收敛在 $\mathbb{R}^{m}\times[0,T_{\omega}]$ 的紧子集上一致。此外，对任意固定 $t\in[0,T_{\omega}]$，$U_{k}(\cdot,t)\rightharpoonup U(\cdot,t)$ 在 $H_{\ell}$ 中弱收敛。

接下来我们证明对任意 $t\in[0,T_{\omega}]$，$U_{k}(\cdot,t)\to U(\cdot,t)$ 在 $H_{0}$ 中强收敛。为此置 $V_{k}=U_{k}-U$，诉诸\eqref{align:ch05-Ul1-est}，它在当前设定下仍成立。然而我们不再能使用\eqref{align:ch05-Vk-positive}，因为 $A(U)$ 不必正定。取而代之地，我们采用
\begin{align}
\int_{\mathbb{R}^{m}}V_{k}^{\top}(x,t)A(U(x,t))V_{k}(x,t)\,\mathrm{d}x\geqslant\mu\|V_{k}(\cdot,t)\|_{0}^{2}-\delta\|V_{k}(\cdot,t)\|_{-1}^{2},\label{align:ch05-Vk-cones}
\end{align}
它来自\cref{lemma:ch05-compactness}。我们可通过注意到 $U$ 满足\eqref{align:ch05-cauchy-system}而 $U_{k}$ 满足\eqref{align:ch05-parabolic}、二者初值相同，来处理项 $\|V_{k}(\cdot,t)\|_{-1}$：
\begin{align}
V_{k}(x,t)=\sum_{\alpha=1}^{m}\partial_{\alpha}\int_{0}^{t}[G_{\alpha}(U(x,\tau))-G_{\alpha}(U_{k}(x,\tau))+\varepsilon_{k}\partial_{\alpha}U_{k}(x,\tau)]\,\mathrm{d}\tau,\label{align:ch05-Vk-rep}
\end{align}
从而
\begin{align}
\|V_{k}(\cdot,t)\|_{-1}\leqslant c\int_{0}^{t}[\|V_{k}(\cdot,\tau)\|_{0}+\varepsilon\|U_{k}(\cdot,\tau)\|_{1}]\,\mathrm{d}\tau.\label{align:ch05-Vk-negnorm}
\end{align}
合并\eqref{align:ch05-Ul1-est}、\eqref{align:ch05-Vk-cones}与\eqref{align:ch05-Vk-negnorm}，我们推得
\begin{align}
\|V_{k}(\cdot,t)\|_{0}^{2}\leqslant c(\omega+t)\int_{0}^{t}\|V_{k}(\cdot,\tau)\|_{0}^{2}\,\mathrm{d}\tau+c\omega^{2}\varepsilon_{k}^{2}t(1+\varepsilon_{k}t),\label{align:ch05-Vk-Gronwall-cones}
\end{align}
它可与\eqref{align:ch05-Vk-Gronwall}比较，并蕴涵对每个 $t\in[0,T_{\omega}]$，$V_{k}(\cdot,t)\to0$ 在 $H_{0}$ 中强收敛。由于 $U_{k}(\cdot,t)\rightharpoonup U(\cdot,t)$ 在 $H_{\ell}$ 中弱收敛，这就给出 $U_{k}(\cdot,t)\to U(\cdot,t)$ 在 $H_{\ell-1}$ 中强收敛且在 $\mathbb{R}^{m}$ 上一致。此外，$\nabla U_{k}(\cdot,t)\to\nabla U(\cdot,t)$ 在 $\mathbb{R}^{m}$ 上一致，对任意 $t\in[0,T_{\omega}]$。特别地，由于 $U_{k}$ 满足\eqref{align:ch05-uniform-est}、$g$ 通过\eqref{align:ch05-g-cones}有界，我们推得
\begin{align}
\|U(\cdot,t)\|_{\ell}\leqslant c_{0}\|U_{0}\|_{\ell}\exp\int_{0}^{t}c[1+\|\|\nabla U(\cdot,\tau)\|\|_{L^{\infty}}]\,\mathrm{d}\tau.\label{align:ch05-cones-est}
\end{align}

下一步是验证 $t\mapsto U(\cdot,t)$ 在 $[0,T_{\omega}]$ 上连续。如\cref{lemma:ch05-limit}证明中那样，考虑恒等式\eqref{align:ch05-continuity-id}，它在当前设定下仍有效，令 $k\to\infty$。\eqref{align:ch05-continuity-id}右端前两项趋于零，而第三项的上极限非正。事实上，即使我们不再有\eqref{align:ch05-continuity-sign}，取而代之，由\cref{lemma:ch05-compactness}，
\begin{align}
\sum_{|r|\leqslant\ell}\int_{\mathbb{R}^{m}}(U_{kr}-U_{r})^{\top}A(U)(U_{kr}-U_{r})\,\mathrm{d}x\geqslant\mu\|U_{k}-U\|_{\ell}^{2}-c\|U_{k}-U\|_{\ell-1}^{2},\label{align:ch05-cones-sign}
\end{align}
且如上所注 $\|U_{k}-U\|_{\ell-1}\to0$。由于\eqref{align:ch05-continuity-id}左端的上极限非正，且 $U_{k}$ 满足\eqref{align:ch05-sum-energy}、$g$ 通过\eqref{align:ch05-g-cones}有界，我们导出不等式
\begin{align}
\sum_{|r|\leqslant\ell}\int_{\mathbb{R}^{m}}U_{r}^{\top}(x,t)A(U(x,t))U_{r}(x,t)\,\mathrm{d}x-\sum_{|r|\leqslant\ell}\int_{\mathbb{R}^{m}}U_{0r}^{\top}(x)A(U_{0}(x))U_{0r}(x)\,\mathrm{d}x\leqslant c\omega^{2}(1+\omega)t,\label{align:ch05-cones-est2}
\end{align}
它是\eqref{align:ch05-continuity-est}的类比。

沿着\cref{lemma:ch05-limit}证明的路线图继续，我们观察到恒等式\eqref{align:ch05-continuity-id2}在当前设定下仍然有效。我们不再有\eqref{align:ch05-continuity-positive}，取而代之，由\cref{lemma:ch05-compactness}，
\begin{align}
\sum_{|r|\leqslant\ell}\int_{\mathbb{R}^{m}}[U_{r}(x,t)-U_{0r}(x)]^{\top}A(U_{0}(x))[U_{r}(x,t)-U_{0r}(x)]\,\mathrm{d}x\geqslant\mu\|U(\cdot,t)-U_{0}(\cdot)\|_{\ell}^{2}-c\|U(\cdot,t)-U_{0}(\cdot)\|_{\ell-1}^{2}.\label{align:ch05-cones-positive}
\end{align}
由\eqref{align:ch05-dtU-bound}，取 $\varepsilon=0$，$\|U(\cdot,t)-U_{0}(\cdot)\|_{\ell-1}=O(t)$，当 $t\to0$。另一方面，作为 $U(\cdot,t)\to U_{0}(\cdot)$ 在 $H_{\ell}$ 中弱收敛连同\eqref{align:ch05-cones-est2}与\eqref{align:ch05-Linf-bound}的推论，当 $t\to0$ 时\eqref{align:ch05-continuity-id2}右端的上极限非正。于是由此得：当 $t\to0$ 时 $\|U(\cdot,t)-U_{0}(\cdot)\|_{\ell}\to0$，即 $t\mapsto U(\cdot,t)$ 在 $t=0$ 处右连续，从而在 $[0,T_{\omega}]$ 上连续。

定理证明的剩余步骤就是\cref{theorem:ch05-凸熵局部存在}证明的最后步骤：将上述构造的\eqref{align:ch05-cauchy-system}、\eqref{align:ch05-cauchy-data}的解 $U$ 延拓到 $T_{\omega}$ 之后，
以初值 $U(\cdot,T_{\omega})$ 解一个新的\eqref{align:ch05-cauchy-system} Cauchy 问题，迭代此过程直到达到极大时间区间 $[0,T_{\infty})$。寿命 $T_{\infty}$ 仅当 $\limsup_{t\to T_{\infty}}\|U(\cdot,t)\|_{L^{\infty}}=\rho$ 且/或 $\limsup_{t\to T_{\infty}}\|U(\cdot,t)\|_{\ell}=\infty$ 时才可能有限。由\eqref{align:ch05-cones-est}，$\|U(\cdot,t)\|_{\ell}$ 仅当\eqref{align:ch05-involution-blowup1}成立时才可能无界。最后，由 $U(\cdot,t)\in C^{0}([0,T_{\infty});H_{\ell})$ 与\eqref{align:ch05-cauchy-system}得\eqref{align:ch05-involution-regularity}。定理证明完成。
\end{proof}

\begin{proof}
\emph{\cref{theorem:ch05-involution-stability}的证明。}只需回溯\cref{theorem:ch05-相对熵稳定性}证明的步骤，作必要调整。

$h(U,\bar{U})$ 与 $Z_{\alpha}(U,\bar{U})$ 的定义\eqref{align:ch05-h-def}与\eqref{align:ch05-Zalpha-def}保持不变。然而，为顾及 $\eta$ 现仅为偶然熵这一假设，$Y_{\alpha}(U,\bar{U})$ 的定义\eqref{align:ch05-Yalpha-def}须替换为
\begin{align}
Y_{\alpha}(U,\bar{U})=Q_{\alpha}(U)-Q_{\alpha}(\bar{U})-\mathrm{D}\eta(\bar{U})[G_{\alpha}(U)-G_{\alpha}(\bar{U})]-\Xi(\bar{U})^{\top}M_{\alpha}[U-\bar{U}].\label{align:ch05-Yalpha-contingent}
\end{align}
作此改变后，取代\eqref{align:ch05-relative-ineq1}，我们现已有
\begin{align}
\int_{0}^{T}\int_{\mathbb{R}^{m}}[\partial_{t}\psi h(U,\bar{U})+\sum_{\alpha=1}^{m}\partial_{\alpha}\psi Y_{\alpha}(U,\bar{U})]\,\mathrm{d}x\,\mathrm{d}t+\int_{\mathbb{R}^{m}}\psi(x,0)h(U_{0}(x),\bar{U}_{0}(x))\,\mathrm{d}x
&\geqslant-\int_{0}^{T}\int_{\mathbb{R}^{m}}\partial_{t}\psi\mathrm{D}\eta(\bar{U})[U-\bar{U}]\,\mathrm{d}x\,\mathrm{d}t\nonumber\\
&\quad-\int_{0}^{T}\int_{\mathbb{R}^{m}}\sum_{\alpha=1}^{m}\partial_{\alpha}\psi\{\mathrm{D}\eta(\bar{U})[G_{\alpha}(U)-G_{\alpha}(\bar{U})]+\Xi(\bar{U})^{\top}M_{\alpha}[U-\bar{U}]\}\,\mathrm{d}x\,\mathrm{d}t.\label{align:ch05-relative-ineq1-contingent}
\end{align}

\eqref{align:ch05-relative-id1}仍然有效，但\eqref{align:ch05-Deta-t}须修改以反映由\eqref{align:ch05-Jalpha-contingent}定义的新 $J_{\alpha}(U)$ 的对称性。顾及 $\bar{U}$ 满足对合\eqref{align:ch05-involution}，我们得
\begin{align}
\partial_{t}\mathrm{D}\eta(\bar{U})=\partial_{t}\bar{U}^{\top}A(\bar{U})&=-\sum_{\alpha=1}^{m}\partial_{\alpha}\bar{U}^{\top}\mathrm{D}G_{\alpha}(\bar{U})^{\top}A(\bar{U})\nonumber\\
&=-\sum_{\alpha=1}^{m}\partial_{\alpha}\bar{U}^{\top}J_{\alpha}(\bar{U})^{\top}=-\sum_{\alpha=1}^{m}\partial_{\alpha}\bar{U}^{\top}J_{\alpha}(\bar{U})\nonumber\\
&=-\sum_{\alpha=1}^{m}\partial_{\alpha}\bar{U}^{\top}A(\bar{U})\mathrm{D}G_{\alpha}(\bar{U})-\sum_{\alpha=1}^{m}\partial_{\alpha}\Xi(\bar{U})^{\top}M_{\alpha}.\label{align:ch05-Deta-t-contingent}
\end{align}

合并\eqref{align:ch05-relative-ineq1-contingent}、\eqref{align:ch05-relative-id1}与\eqref{align:ch05-Deta-t-contingent}，并利用 $U-\bar{U}$ 满足对合\eqref{align:ch05-involution}，我们结论\eqref{align:ch05-relative-ineq2}在当前设定下仍然有效。我们固定 $L^{\infty}$ 弱$^{*}$ 连续性 $\eta(U(\cdot,t))$ 的任意点 $\sigma\in(0,T)$ 与充分小的 $\varepsilon$，对试验函数 $\psi(x,t)=\chi(x)\omega(t)$ 应用\eqref{align:ch05-relative-ineq2}，其中 $\omega$ 由\eqref{align:ch05-cutoff-t}定义，$\chi$ 是任意 $C_{0}^{\infty}(\mathbb{R}^{m})$ 函数，对所有 $|x|<\varepsilon^{-1}$ 的 $x$ 有 $\chi(x)=1$。令 $\varepsilon\to0$，我们导出
\begin{align}
\int_{\mathbb{R}^{m}}h(U(x,\sigma),\bar{U}(x,\sigma))\,\mathrm{d}x
\leqslant\int_{\mathbb{R}^{m}}h(U_{0}(x),\bar{U}_{0}(x))\,\mathrm{d}x-\int_{0}^{\sigma}\int_{\mathbb{R}^{m}}\sum_{\alpha=1}^{m}\partial_{\alpha}\bar{U}^{\top}Z_{\alpha}(U,\bar{U})\,\mathrm{d}x\,\mathrm{d}t.\label{align:ch05-cones-integral}
\end{align}

由\eqref{align:ch05-h-def}，
\begin{align}
h(U,\bar{U})=\frac{1}{2}(U-\bar{U})^{\top}H(U,\bar{U})(U-\bar{U}),\label{align:ch05-h-Hessian}
\end{align}
其中
\begin{align}
H(U,\bar{U})=2\int_{0}^{1}\int_{0}^{z}A(sU+(1-s)\bar{U})\,\mathrm{d}s\,\mathrm{d}z.\label{align:ch05-H-def}
\end{align}
由\eqref{align:ch05-cone-convex}，
\begin{align}
X^{\top}H(U,\bar{U})X\geqslant 2\mu|X|^{2},\qquad X\in\mathscr{C},\ U\in\overline{\mathscr{B}_{\rho}},\ \bar{U}\in\overline{\mathscr{B}_{\rho}}.\label{align:ch05-H-positive}
\end{align}

由于 $\bar{U}$ Lipschitz，当\eqref{align:ch05-small-oscillation}以充分小的 $\kappa$ 成立时，存在 $\varepsilon>0$ 使得 $|x-y|<2\varepsilon$ 蕴涵
\begin{align}
|H(U(x,t),\bar{U}(x,t))-H(U(y,t),\bar{U}(y,t))|<\frac{1}{2}\mu,\label{align:ch05-H-oscillation}
\end{align}
对所有 $t\in[0,T)$ 成立。此外，顾及 $U(x,t)$ 与 $\bar{U}(x,t)$ 当 $|x|\to\infty$ 的规定行为，对大 $\alpha$，\eqref{align:ch05-H-oscillation} 对所有 $t\in[0,T)$ 及集合 $\Omega_{0}=\{z\in\mathbb{R}^{m}:|z|>\alpha\}$ 中所有 $x,y$ 亦成立。我们以半径 $\varepsilon$ 的球 $\Omega_{1},\ldots,\Omega_{K}$ 固定紧集 $\Omega_{0}^{c}$ 的某个覆盖。于是，对任意固定 $t\in[0,T)$，$\Omega_{0},\Omega_{1},\ldots,\Omega_{K}$ 给出 $\mathbb{R}^{m}$ 的一个覆盖，满足\cref{lemma:ch05-compactness}中以 $H(U(\cdot,t),\bar{U}(\cdot,t))$ 扮演 $P(\cdot)$ 的条件。因此
\begin{align}
\int_{\mathbb{R}^{m}}h(U(x,\sigma),\bar{U}(x,\sigma))\,\mathrm{d}x\geqslant\frac{1}{2}\mu\|U(\cdot,\sigma)-\bar{U}(\cdot,\sigma)\|_{0}^{2}-\frac{1}{2}\delta\|U(\cdot,\sigma)-\bar{U}(\cdot,\sigma)\|_{-1}^{2}.\label{align:ch05-h-lower}
\end{align}
由
\begin{align}
U(\cdot,\sigma)-\bar{U}(\cdot,\sigma)=U_{0}(\cdot)-\bar{U}_{0}(\cdot)-\sum_{\alpha=1}^{m}\partial_{\alpha}\int_{0}^{\sigma}[G_{\alpha}(U(\cdot,t))-G_{\alpha}(\bar{U}(\cdot,t))]\,\mathrm{d}t,\label{align:ch05-difference-rep}
\end{align}
得
\begin{align}
\|U(\cdot,\sigma)-\bar{U}(\cdot,\sigma)\|_{-1}\leqslant\|U_{0}(\cdot)-\bar{U}_{0}(\cdot)\|_{0}+c\int_{0}^{\sigma}\|U(\cdot,t)-\bar{U}(\cdot,t)\|_{0}\,\mathrm{d}t.\label{align:ch05-difference-negnorm}
\end{align}

合并\eqref{align:ch05-cones-integral}、\eqref{align:ch05-h-lower}、\eqref{align:ch05-difference-negnorm}与\eqref{align:ch05-Zalpha-def}，我们推得由 $u(t)=\|U(\cdot,t)-\bar{U}(\cdot,t)\|_{0}$ 定义的函数 $u$ 满足积分不等式
\begin{align}
u(\sigma)\leqslant au(0)+b\int_{0}^{\sigma}u(t)\,\mathrm{d}t,\label{align:ch05-cones-Gronwall}
\end{align}
对几乎所有 $\sigma\in[0,T)$ 成立，其中 $a$ 仅依赖于 $G,\eta$ 及其在 $\overline{\mathscr{B}_{\rho}}$ 上导数的界，而 $b$ 还依赖于 $T$ 与 $\bar{U}$ 的 Lipschitz 常数。对\eqref{align:ch05-cones-Gronwall}应用 Gronwall 不等式，即得\eqref{align:ch05-involution-stability-est}。证明完成。
\end{proof}

还应注意，在某些守恒律方程组\eqref{align:ch05-cauchy-system}中，通过初值引入的形如\eqref{align:ch05-involution}的关系，即使 $M_{\alpha}$ 并非由对称关系\eqref{align:ch05-symmetry-condition}诱导，也可能被（至少经典）解保持。一个典型例子是无旋流条件 $\omega=\operatorname{curl}v=0$，它被 Euler 方程组\eqref{align:ch03-协调跳跃分解}的光滑解保持，
但在间断发展后失效。

还有装备有涉及状态向量及其空间导数的非线性函数的对合的守恒律方程组。例如，\eqref{align:ch02-2-2-12}与\eqref{align:ch02-2-2-13}

可视为等熵弹性动力学方程组\eqref{align:ch03-等熵热弹性系统}的非线性对合。
此外，若在 Euler 坐标中写出\eqref{align:ch03-等熵热弹性系统}，则线性对合\eqref{align:ch03-协调方程}变为非线性：

$F_{j\beta}\partial_{j}F_{i\alpha}-F_{j\alpha}\partial_{j}F_{i\beta}=0$（服从求和约定）。然而，带非线性对合的方程组最著名的例子由广义相对论的 Einstein 方程提供。

\cref{theorem:ch05-involution-existence}、\cref{theorem:ch05-involution-stability}及其证明证实：从分析的角度看，偶然熵概念是装备有对合的守恒律方程组熵概念的自然推广。然而，在上述定理对连续介质物理中出现的方程组的应用中，如\eqref{align:ch03-等熵热弹性系统}，
$\eta$ 是真正的熵。非熵的偶然熵的重要性将在下一节阐明。

应用\cref{theorem:ch05-involution-existence}与\cref{theorem:ch05-involution-stability}的前提并非物理中遇到的每个带非凸熵与对合的方程组都满足。例如，如上所述，在电动力学方程组\eqref{align:ch03-Maxwell}中（$J=0$），
装备有对合\eqref{align:ch03-高斯}（$\rho=0$），
对合锥是整个状态空间 $\mathbb{R}^{6}$，电磁能 $\eta$ 在其上不凸。即使当应用这些定理的条件具备时，如等熵弹性动力学方程组\eqref{align:ch03-等熵热弹性系统}的情形，
也应注意，与\cref{theorem:ch05-相对熵稳定性}相比，\cref{theorem:ch05-involution-stability}要求更多而给出更少。
下一节我们将识别与连续介质物理中遇到的某些方程组中偶然熵的存在相联系的特殊结构，它可弥补上述不足。

\end{document}
